% !TeX root = ../main.tex
\section{Block-Size Sweep}\label{sec:block-size}
\rubric{Block-size and pattern sweeps}{35}{Complete matrix and clear interpretation of IOPS and MiB/s
against latency.}

\begin{hypothesis}
  As the block grows from \qty{4}{\kibi\byte} to \qty{256}{\kibi\byte} at QD=1, small-block I/O should be
  \emph{IOPS-limited} (throughput $=$ IOPS $\times$ block, rising with block size) and large-block I/O
  should approach a \emph{throughput ceiling} where IOPS falls as $1/\text{block}$. Sequential should beat
  random at every block size, with the gap largest at small blocks where per-op overhead dominates.
\end{hypothesis}

\begin{knobs}
  Fixed       & iodepth 1; one job; \code{O_DIRECT}; 100\% read \\
  Varied      & block size \qtyrange{4}{256}{\kibi\byte}; pattern random vs.\ sequential \\
  Script      & \code{SSD_Performance_Profiling/scripts/run.sh block_size} \\
\end{knobs}

\begin{figure}[tbp]
  \centering
  \IfDataExists{block_size}{%
    \begin{subfigure}[t]{0.49\linewidth}
      \begin{tikzpicture}
        \begin{axis}[acs, height=0.85\linewidth, xmode=log, log basis x=2,
                     xlabel={Block size (bytes)}, ylabel={Throughput (MiB/s)}, ymin=0, legend pos=north west]
          \addplot+[iqr bars] table[x=bs, y=mibps_median, iqr of=mibps, only if={pattern}{rand}] {data/block_size.csv};
          \addlegendentry{random}
          \addplot+[iqr bars] table[x=bs, y=mibps_median, iqr of=mibps, only if={pattern}{seq}] {data/block_size.csv};
          \addlegendentry{sequential}
        \end{axis}
      \end{tikzpicture}
      \caption{Throughput.}
    \end{subfigure}\hfill
    \begin{subfigure}[t]{0.49\linewidth}
      \begin{tikzpicture}
        \begin{axis}[acs, height=0.85\linewidth, xmode=log, log basis x=2, ymode=log,
                     xlabel={Block size (bytes)}, ylabel={Mean latency (\si{\micro\second})}]
          \addplot+[iqr bars] table[x=bs, y=lat_mean_us_median, iqr of=lat_mean_us, only if={pattern}{rand}] {data/block_size.csv};
          \addplot+[iqr bars] table[x=bs, y=lat_mean_us_median, iqr of=lat_mean_us, only if={pattern}{seq}] {data/block_size.csv};
        \end{axis}
      \end{tikzpicture}
      \caption{Mean latency (same runs).}
    \end{subfigure}%
  }
  \caption{QD=1 read block-size sweep. Random throughput rises with block size while IOPS stays bounded;
  sequential saturates near \qty{1.5}{\gibi\byte\per\second}. Random is far slower at small blocks and the
  gap narrows as transfers dominate per-op cost.}\label{fig:block-size}
\end{figure}

\subsection{Analysis}
The two patterns behave exactly as the IOPS-vs-throughput model predicts. \textbf{Sequential} reads start
at \num{87}k IOPS / \qty{341}{\mebi\byte\per\second} for \qty{4}{\kibi\byte} and climb to a throughput
plateau of \ensuremath{\sim}\qty{1.5}{\gibi\byte\per\second} by \qty{256}{\kibi\byte} (IOPS then falling as $1/$block:
\num{6013} IOPS at \qty{256}{\kibi\byte}); the device is throughput-limited once the block is large enough
to amortise per-command overhead. \textbf{Random} reads are IOPS-bound: IOPS only drops from \num{20240}
(\qty{4}{\kibi\byte}) to \num{3885} (\qty{256}{\kibi\byte}) while throughput rises from
\qty{79}{\mebi\byte\per\second} to \qty{971}{\mebi\byte\per\second}, because each QD=1 random access pays
the fixed \ensuremath{\sim}\qty{50}{\micro\second} flash-read latency regardless of size (mean latency grows only
\qty{49}{}$\to$\qty{257}{\micro\second} over a $64\times$ block range). The sequential/random gap is
widest at \qty{4}{\kibi\byte} ($4.3\times$ throughput) and narrows to $1.5\times$ at
\qty{256}{\kibi\byte}, where both approach the same per-transfer regime. Note these are QD=1 numbers; the
device's full sequential bandwidth would need a deeper queue or larger blocks, so \qty{1.5}{\gibi\byte\per\second}
is a single-request floor, not the interface ceiling.
